\section{Sensitivity and limits}

The broad result depends on exact-option continuation scoring under canonical answer ordering. That is the intended primary readout because it scores the answer string itself rather than a proxy label. Alternate readouts and option orders are useful sensitivity material, but they can change the final winner and therefore are not used to replace the primary path.

Several edge cases are handled explicitly. Final-layer ties are logged, and rows with nonpositive final gaps are excluded from normalized-gap metrics rather than forced into $d_{50}$. Across included broad-panel rows, \FinalLayerTieCount{} rows have a final-layer top tie and \NormalizedGapExcludedCount{} rows are excluded from normalized-gap metrics. The appendix reports neighboring thresholds, $25\%$ and $75\%$, to show how the margin-depth timing shifts when the criterion is relaxed or tightened.

The mechanistic results also have a boundary. The attention/MLP decomposition is a scalar accounting model over cached readouts. Span deletion is a true input-level intervention in the legacy cache. Simulated ablation and substitution are linearized bookkeeping estimates: they do not let later layers react to a modified activation state. We therefore use them to support conditional component-level patterns, not to claim complete causal circuits.

The empirical scope is multiple-choice cached readouts in the tested open-weight models and benchmarks. The paper does not claim behavior for open-ended generation, all architectures, or all answer formats.
